\uuid{9jTu}
\titre{Résolution d'équation différentielle du premier ordre}
\auteur{bodin ; croizat ; sacre}
\organisation{exo7}
\licence{CC BY-NC-SA 4.0}
\datecreate{2022-03-08}
\niveau{L1}
\module{Analyse}
\chapitre{Equation différentielle}
\sousChapitre{Résolution d'équation différentielle du premier ordre}
\difficulte{1}
\theme{}
\exo7qcmid{hilisit.equations-differentielles.041}

\contenu{
\begin{question}
\qclassification{225.01}
Soit l'équation différentielle $y' = 2y + \cos(x)$.
Quelles sont les affirmations vraies ?
\begin{answers}
  \bad{Les solutions de l'équation homogène associée sont les $y(x) = C\sin(x)$.}
  \bad{Les solutions de l'équation homogène associée sont les $y(x) = C\cos(x)$.}
  \good{Une solution particulière est $y(x) = \frac15\sin(x)-\frac25\cos(x)$.}
  \bad{Une solution particulière est $y(x) = e^{2x}$.}
\end{answers}
\begin{explanations}
L'équation homogène est $y'=2y$, dont les solutions sont les $y_h(x) = Ce^{2x}$.
Une solution particulière de l'équation $y' = 2y +\cos(x)$ est $y_p(x) = \frac15\sin(x)-\frac25\cos(x)$.
Les solutions générales sont alors $y(x) = y_h(x) + y_p(x)$.
\end{explanations}
\end{question}
}
